% 日本語版 CV（assets/pdf/cv_ja.pdf の原稿）
% ビルド: lualatex -output-directory=../assets/pdf cv_ja.tex
\documentclass[a4paper,10pt]{article}

\usepackage[haranoaji,deluxe]{luatexja-preset}
\usepackage{fontspec}
\setmainfont{XCharter}[
  Extension=.otf, UprightFont=*-Roman, BoldFont=*-Bold,
  ItalicFont=*-Italic, BoldItalicFont=*-BoldItalic]
\setsansfont{TeX Gyre Heros}

\usepackage[top=16mm,bottom=18mm,left=17mm,right=17mm]{geometry}
\usepackage{xcolor}
\usepackage{titlesec}
\usepackage{enumitem}
\usepackage{fancyhdr}
\usepackage{needspace}
\usepackage{etoolbox}
\usepackage[hidelinks,pdfusetitle]{hyperref}

\definecolor{accent}{HTML}{1F3A5F}
\definecolor{muted}{HTML}{6B7A90}
\definecolor{rulec}{HTML}{8EA3BF}

\hypersetup{
  pdftitle={許 以寧（Yining Xu）- 履歴書（CV）},
  pdfauthor={許 以寧（Yining Xu）},
}

\setlength{\parindent}{0pt}
\setlength{\parskip}{0pt}
\ltjsetparameter{jacharrange={-2}}
\linespread{1.08}

% セクション見出し
\titleformat{\section}
  {\color{accent}\gtfamily\bfseries\large}
  {}{0pt}{}[\vspace{-0.6ex}{\color{rulec}\titlerule[0.6pt]}]
\titlespacing*{\section}{0pt}{2.2ex}{1.0ex}
\titleformat{\subsection}
  {\color{accent}\gtfamily\bfseries\normalsize}{}{0pt}{}
\titlespacing*{\subsection}{0pt}{1.4ex}{0.5ex}
\titleformat{\subsubsection}
  {\gtfamily\bfseries\small}{}{0pt}{}
\titlespacing*{\subsubsection}{0pt}{0.8ex}{0.3ex}

% 見出しがページ末尾に孤立しないようにする
\pretocmd{\section}{\needspace{6\baselineskip}}{}{}
\pretocmd{\subsection}{\needspace{4\baselineskip}}{}{}

% 箇条書き
\setlist[itemize]{leftmargin=1.2em,itemsep=0.1ex,topsep=0.3ex,parsep=0pt,
  label={\color{muted}\textbullet}}
\setlist[enumerate]{leftmargin=1.8em,itemsep=0.4ex,topsep=0.3ex,parsep=0pt}

% 項目: \entry{見出し}{期間}{所属}
\newcommand{\entry}[3]{%
  \par\addvspace{0.9ex}%
  {\gtfamily\bfseries #1}\hfill{\small\color{muted}#2}\par
  \ifx\relax#3\relax\else{\small #3}\par\fi}
% 行末に右寄せで期間を置く（入らなければ次行の右端へ）
\newcommand{\when}[1]{{\unskip\nobreak\hfil\penalty50\hskip1.5em\hbox{}\nobreak\hfill
  \mbox{\small\color{muted}#1}\parfillskip=0pt\finalhyphendemerits=0\par}}
\newcommand{\me}{\textbf{Xu YS}}

% フッター
\pagestyle{fancy}
\fancyhf{}
\renewcommand{\headrulewidth}{0pt}
\fancyfoot[R]{\footnotesize\color{muted}許 以寧 · 履歴書（CV）· \thepage}

\begin{document}

% ===== ヘッダー =====
\begin{center}
  {\color{accent}\gtfamily\bfseries\fontsize{22}{26}\selectfont 許\, 以寧}\\[0.5ex]
  {\color{accent}\sffamily\bfseries Yining Xu}\\[0.8ex]
  {\small 京都大学大学院 総合生存学館 博士課程}\\[0.3ex]
  {\small\color{muted}〒606-8501 京都府京都市左京区吉田本町}\\[0.4ex]
  {\small
    \href{mailto:ineikyos@gmail.com}{ineikyos@gmail.com}\quad{\color{rulec}|}\quad
    \href{https://yiningsxu.github.io/}{yiningsxu.github.io}\quad{\color{rulec}|}\quad
    \href{https://github.com/yiningsxu}{github.com/yiningsxu}\quad{\color{rulec}|}\quad
    \href{https://researchmap.jp/inei_yining}{researchmap.jp/inei\_yining}}
\end{center}
\vspace{-1ex}

% ===== 研究関心 =====
\section{研究関心}
公衆衛生学とデータサイエンスを専門とし、児童虐待予防、母子保健、口腔保健に重点を置いている。多機関の行政・臨床データを統合し、統計モデリングと解釈可能な機械学習により、早期のリスク検出とエビデンスに基づく政策立案を支援する。また、非線形時系列解析や経験的動的モデリングを用いた感染症疫学研究にも取り組んでいる。

\vspace{0.6ex}
{\gtfamily\bfseries キーワード:} 児童虐待 · デンタルネグレクト · 母子保健 · 感染症疫学 · 機械学習 · 時系列解析 · 経験的動的モデリング（EDM）

% ===== 学歴 =====
\section{学歴}
\entry{博士課程（総合学術）在籍}{2024年4月 – 現在}{京都大学大学院 総合生存学館（5年一貫制博士課程）}
\begin{itemize}
  \item 博士論文：\textit{歯科健診・生活習慣データを用いた児童虐待の早期リスク推定：モデル開発と実装への道筋}
  \item 一時保護児約2,480名の歯科健診・生活習慣の統合データを用い、多項ロジスティック回帰による解釈可能なスクリーニングモデルを開発中。
\end{itemize}

\entry{修士相当（Qualifying Examination 合格）}{2022年4月 – 2024年3月}{京都大学大学院 総合生存学館}
\begin{itemize}
  \item 論文：\textit{COVID-19パンデミックが日本の出生率に与えた影響の要因分析：流行波と経済指標を用いた時系列・都道府県別相関分析（2015–2022）}
  \item 主要な流行波の約9か月後に出生数が有意に減少し、都道府県の一人当たり所得の低下が出生数の減少と関連することを示した。
\end{itemize}

\entry{学士（理学）}{2018年4月 – 2022年3月}{奈良女子大学 理学部 数物科学科 数学コース}
\begin{itemize}
  \item 位相幾何学・代数学・解析学を履修。ディープラーニング実装および Bishop『パターン認識と機械学習』の輪読に参加。
\end{itemize}

% ===== 研究歴 =====
\section{研究歴}
\entry{奨励研究員（JST SPRING）}{2025年4月 – 現在}{京都大学 大学院教育支援機構}

\entry{特別研究員}{2025年4月 – 現在}{帝京大学大学院 公衆衛生学研究科}

\entry{リサーチインターン}{2026年8月 – 2026年9月}{日本電気株式会社（NEC）研究開発部門}
\begin{itemize}
  \item 医療分野における大規模言語モデルの応用に関する研究に従事。
\end{itemize}

\entry{リサーチアシスタント}{2024年4月 – 2025年3月}{京都大学}
\begin{itemize}
  \item ノンパラメトリック・非線形時系列手法を用いたデータ駆動型の感染症モデルを開発。
  \item EDMに基づく因果推論により、感染症間の相互作用（干渉、同時流行）を推定。
\end{itemize}

\entry{リサーチアシスタント}{2023年4月 – 2024年3月}{京都大学}
\begin{itemize}
  \item 感染症動態の数理解析を行い、疫学研究のデータ処理・モデリングを支援。
\end{itemize}

\entry{リサーチアシスタント}{2022年9月 – 2023年2月}{京都大学}
\begin{itemize}
  \item 高齢者の余暇活動が脳の構造・機能に与える影響に関する縦断研究に従事。
  \item 尿・唾液サンプル採取やコルチゾール関連手順を含む、生化学・MRI関連の研究室業務を補助。
\end{itemize}

% ===== 論文 =====
\section{論文}
{\small\color{muted}* 責任著者}

\subsection{児童虐待・口腔保健}
\subsubsection{査読付き論文}
\begin{enumerate}
  \item \me*, Kohashi K, Morioka S, Mizumoto K. Oral health differences across maltreatment categories among children in temporary protection facilities in Japan: a retrospective cross-sectional study. \textit{BMC Public Health}. 2026.（印刷中）
\end{enumerate}
\subsubsection{執筆中}
\begin{enumerate}
  \item \me, et al. Prevalence of early childhood caries and associated factors among preschool children in temporary protective custody. 投稿予定: \textit{International Journal of Paediatric Dentistry}.
\end{enumerate}

\subsection{感染症疫学}
\subsubsection{査読付き論文}
\begin{enumerate}
  \item \me, Takahashi K, Shimakawa Y, Nakamura T, Itokazu T, Takayama Y, Mizumoto K. Timing of outbreak-triggered unit-wide screening and COVID-19 outbreak size in long-term care and related social welfare facilities in Okinawa, Japan. \textit{IJID Regions}. 2026;101042. \href{https://doi.org/10.1016/j.ijregi.2026.101042}{doi:10.1016/j.ijregi.2026.101042}
  \item Kozuka M, \me, Chowell G, Ogawa K, Takayama Y, Mizumoto K. Traveler-derived importation risk underestimates regional disease activity: evidence from Okinawa, Japan. \textit{Journal of Infection and Public Health}. 2026;103298. \href{https://doi.org/10.1016/j.jiph.2026.103298}{doi:10.1016/j.jiph.2026.103298}
  \item Ogawa K, \me, Shimakawa Y, Chowell G, Omori R, Matsuyama R, Takayama Y, Mizumoto K. Secondary severe acute respiratory syndrome coronavirus 2 transmission from childcare workers versus teachers in school-associated screening events, Okinawa, Japan, January–March 2022. \textit{International Journal of Infectious Diseases}. 2026;167:108672. \href{https://doi.org/10.1016/j.ijid.2026.108672}{doi:10.1016/j.ijid.2026.108672}
  \item \me, Shimakawa Y, Chowell G, Matsuyama R, Nagamoto T, Omori R, Nakamura T, Itokazu T, Takayama Y, Mizumoto K. Determinants of COVID-19 outbreak size in elderly residential facilities in Okinawa Prefecture, Japan, April to June 2022. \textit{IJID Regions}. 2025;100813. \href{https://doi.org/10.1016/j.ijregi.2025.100813}{doi:10.1016/j.ijregi.2025.100813}
  \item Kozuka M, Shimakawa Y, Chowell G, Nagamoto T, Matsuyama R, Omori R, \me, Shimoji Y, Ogawa K, Takayama Y, Mizumoto K. Pre-travel health awareness and perceptions of voluntary airport PCR testing during COVID-19: a cross-sectional study in Okinawa, Japan. \textit{IJID Regions}. 2025;100817. \href{https://doi.org/10.1016/j.ijregi.2025.100817}{doi:10.1016/j.ijregi.2025.100817}
  \item Takayama Y, \me, Shimakawa Y, Chowell G, Kozuka M, Omori R, Matsuyama R, Yamamoto T, Mizumoto K. Assessment of fever screening at airports in detecting domestic passengers infected with SARS-CoV-2, 2020–2022, Okinawa Prefecture, Japan. \textit{BMC Infectious Diseases}. 2024;24(1):542. \href{https://doi.org/10.1186/s12879-024-09427-5}{doi:10.1186/s12879-024-09427-5}
\end{enumerate}
\subsubsection{査読中}
\begin{enumerate}
  \item Tu S, \me, et al. Information sources, vaccine literacy, and intention to receive a COVID-19 booster dose: a cross-sectional online survey in Okinawa, Japan. \textit{Vaccine: X}.
\end{enumerate}
\subsubsection{執筆中}
\begin{enumerate}
  \item \me, et al. Implications of subtype interference for seasonal influenza control: an empirical dynamic modeling study in Japan. 投稿予定: \textit{Eurosurveillance}.
\end{enumerate}

% ===== 学会発表 =====
\section{学会発表}
\subsection{児童虐待・口腔保健}
\begin{enumerate}
  \item \me, et al. 一時保護所入所児における発達段階および主たる保護理由別にみた口腔保健ニーズ. 日本子ども虐待防止学会 第32回学術集会おおさか大会; 2026年11月15日.（口頭発表、発表予定）
  \item \me, et al. 一時保護児における乳幼児う蝕の有病状況. 第85回日本公衆衛生学会総会（東京）; 2026年10月29日.（ポスター発表、発表予定）
  \item \me, et al. 一時保護児の口腔内状況から読み解く虐待類型. 日本子ども虐待防止学会 第31回学術集会ほっかいどう大会; 2025年11月.（口頭発表）
  \item （共同演者）被虐待児の口腔内状況に影響する生活習慣の調査 子ども虐待のリスク判定の一助として. 日本子ども虐待防止学会 第31回学術集会ほっかいどう大会; 2025年11月.（口頭発表）
  \item \me, et al. 一時保護された被虐待児童の口腔内状況に関する報告. 第95回日本衛生学会学術総会; 2025年3月.（口頭発表）
  \item \me. 児童相談所の一時保護児における口腔内状況と虐待分類から見えてくること. 第9回日本子ども虐待防止歯科研究会・学術大会; 2024年12月.（招待講演）
\end{enumerate}

\subsection{感染症疫学}
\begin{enumerate}
  \item \me, et al. 人流データを用いた震災後の感染症流入リスクに関する報告（第一報）. 第84回日本公衆衛生学会総会; 2025年10月.（口頭発表）
  \item \me, et al. 日本におけるインフルエンザ亜型の動態：時系列解析からの知見. 第95回日本衛生学会学術総会; 2025年3月.（ポスター発表）
  \item \me, et al. Prevent outbreaks in elderly care facilities: factors associated with COVID-19 outbreak size in elderly care facilities in Okinawa Prefecture, Japan, April to June 2022. Joint Congress on Global Health 2024; 2024年11月.（国際学会、ポスター発表）
  \item \me, et al. Interaction analysis between influenza subtypes in Okinawa using a time series analysis. Joint Congress on Global Health 2024; 2024年11月.（国際学会、ポスター発表）
  \item \me, et al. 沖縄県の高齢者介護施設における新型コロナウイルス感染症発生規模に関連する要因. 第83回日本公衆衛生学会総会; 2024年10月.（口頭発表）
\end{enumerate}

% ===== 研究費 =====
\section{研究費・フェローシップ}
\entry{公益財団法人8020推進財団 令和7年度 8020公募研究事業}{2025年6月 – 2026年3月}{\textit{児童相談所の一時保護児における口腔内状況と虐待分類との関連性の定量化・可視化}}

\entry{JST 次世代研究者挑戦的研究プログラム（SPRING）JPMJSP2110}{2025年4月 – 2027年3月}{\textit{Innovative strategies for child abuse prevention: a collaborative study integrating dentistry, machine learning, and child guidance services}}

% ===== 研究プロジェクト =====
\section{研究プロジェクト}
\subsection{児童虐待・口腔保健}
\entry{一時保護児の口腔内状況の特徴と虐待分類の関連に係る分析調査}{}{京都大学（代表）、A県歯科医師会}
\begin{itemize}
  \item 研究代表者：研究デザイン、データ管理、解析、論文執筆
\end{itemize}

\entry{臨床法医学での虐待評価の後方視的研究}{}{C大学（代表）、京都大学}
\begin{itemize}
  \item 研究分担者：データ解析リード
\end{itemize}

\entry{消防職員・病院救命士・救急救命士課程の学生の児童虐待対応に関する現状と教育的介入の影響に関する研究}{}{T大学（代表）、京都大学}
\begin{itemize}
  \item 研究分担者：質問票設計、データ解析
\end{itemize}

\vspace{0.6ex}
\begin{itemize}
  \item {\gtfamily\bfseries 予定：}子ども虐待分野の学会 研究委員会メンバーとして、全国児童虐待レジストリの解析を担当。
\end{itemize}

\subsection{感染症疫学}
\begin{itemize}
  \item {\gtfamily\bfseries 空港における発熱スクリーニングの評価（沖縄、2020–2022）：}データ可視化
  \item {\gtfamily\bfseries 高齢者施設における感染対策とCOVID-19アウトブレイク規模（沖縄、2022）：}データ解析リード
  \item {\gtfamily\bfseries 経験的動的モデリングによる感染症間相互作用の推定：}データ解析リード
  \item {\gtfamily\bfseries 人流データを用いた震災後の感染症流入リスク：}データ解析リード
\end{itemize}

% ===== 受賞 =====
\section{受賞・奨学金}
\begin{itemize}
  \item ロータリー米山記念奨学会 学部課程奨学金\when{2021年4月 – 2022年3月}
  \item 奈良女子大学 佐保会奨学金（特に優秀な学生として表彰）\when{2021年11月}
  \item 文部科学省外国人留学生学習奨励費（追加給付）\when{2020年10月}
  \item 文部科学省外国人留学生学習奨励費\when{2018年4月 – 2019年3月}
\end{itemize}

% ===== 学会活動 =====
\section{学会活動}
\begin{itemize}
  \item 子ども虐待分野の学会 研究委員会委員（解析担当）
\end{itemize}

% ===== 教育・アウトリーチ =====
\section{教育・アウトリーチ}
\begin{itemize}
  \item {\gtfamily\bfseries ゲスト講師}、英語授業・中国文化紹介、奈良女子大学附属小学校\when{2021年10月・12月}
  \item {\gtfamily\bfseries ティーチングアシスタント}、大阪教育大学附属高等学校天王寺校舎の大学訪問プログラム\when{2019年12月}
  \item {\gtfamily\bfseries 研究グループリーダー}、NWU CORE of STEM International Summer Programme（光学と折り紙の位相幾何、6か国の参加者）\when{2019年8月}
  \item {\gtfamily\bfseries ボランティアチューター}、Glocal Camp 2019（京都）：高校生の面接・英語プレゼンテーションを指導（チームが1位を獲得）\when{2019年3月}
  \item {\gtfamily\bfseries ゲスト講師}、小学生向けに折り紙の展開図の位相幾何に関する40分の講義、Junru Primary School（中国）\when{2019年3月}
  \item {\gtfamily\bfseries 数学・英語講師}（アルバイト）、個別指導塾（奈良）、担当生徒の英検合格率100\%\when{2018年11月 – 2020年12月}
  \item {\gtfamily\bfseries 留学生チューター}、京都大学（2022年4月–9月）、奈良女子大学（2021年4月–2022年3月）\when{2021年 – 2022年}
\end{itemize}

% ===== スキル =====
\section{スキル}
\begin{itemize}
  \item {\gtfamily\bfseries 統計手法：}観察研究デザイン（横断研究、時系列、地域比較）、回帰・多変量モデリング（ロジスティック、多項、負の二項、GLM）、時系列・因果分析（季節調整、EDM）、倫理審査下での行政・臨床センシティブデータの管理
  \item {\gtfamily\bfseries 機械学習：}教師あり分類（ランダムフォレスト）、特徴量設計、交差検証、モデル解釈（SHAP、特徴量重要度）、独立成分分析（ICA）
  \item {\gtfamily\bfseries LLM・自然言語処理：}LLMを用いた情報抽出、長文の非構造化テキストを分析可能な構造化データへ変換
  \item {\gtfamily\bfseries プログラミング：}Python（pandas、NumPy、scikit-learn、statsmodels、投稿用図表の自動生成）、R（GLM、EDA、rEDM、rUIC）、JavaScript/HTML/CSS、Java、C（基礎）
  \item {\gtfamily\bfseries ツール：}Git/GitHub、GitHub Pages（Jekyll）、YAML、VS Code
\end{itemize}

% ===== 言語 =====
\section{言語}
\begin{itemize}
  \item {\gtfamily\bfseries 中国語：}ネイティブ
  \item {\gtfamily\bfseries 日本語：}ネイティブに近いレベル（JLPT N1）
  \item {\gtfamily\bfseries 英語：}ビジネスレベル（IELTS 7.0、TOEIC 855）
\end{itemize}

\end{document}
